% Source: https://artofproblemsolving.com/community/c3h1225669_2011_arml_number_8
% Stable identity: aops-topic-1225669-q08
Write $f(x)=a(x-2)(x-r)$. Then $f(f(x))=0$ exactly when $f(x)=2$ or $f(x)=r$. Each of these equations has $0$, $1$ (double) or $2$ real solutions, symmetric about the axis of the parabola. Since $5$ is the only real zero, the axis is $x=5$, so $\dfrac{2+r}{2}=5$ and $r=8$; and one equation has the double root $5$ while the other has none.

$f(5)=-9a$ is the vertex value. If $f(5)=2$, then $a=-\dfrac{2}{9}$, the parabola opens down with maximum $2$, and $f(x)=8$ has no solution: consistent. If $f(5)=8$, then $f(x)=2<8$ has two solutions: impossible. So $f(x)=-\dfrac{2}{9}(x-2)(x-8)$ and $f(0)=-\dfrac{2}{9}\cdot16=\boxed{-\dfrac{32}{9}}$.
